%! TeX root = ./main.tex
% vim: set spelllang=en:
\section{superrec2 --- Segmental reconciliation algorithms}\label{sec:logiciel:superrec2}
\programdetails{https://github.com/UdeM-LBIT/superrec2}{GPL v3.0}{Between February~2021 and April~2026, 4\,000 lines of code}

The \emph{superrec2} software, implemented in Python, served as a testbed for the development and application of the various segmental reconciliation algorithms presented in this manuscript.
As such, several adjacent problems are solved in it in addition to the implementation of the algorithms themselves:

\begin{compactitemize}
    \item modelling segmental reconciliation scenarios and representing them in textual form to allow the results of the algorithms to be saved;
    \item validating these scenarios in order to experimentally check that the results produced by the algorithms are consistent with the defined models;
    \item visually representing the scenarios to allow their inspection.
\end{compactitemize}

The \textbf{\texttt{superrec2.render}} module and the \textbf{\texttt{superrec2 draw}} command-line program make it possible to load a bilateral or segmental reconciliation scenario and to generate a visual representation of it following the conventions adopted in \cref{fig:reconciliation:rec-trees,fig:reconciliation:rec-dtlx,fig:reconciliation:time-cons,fig:evol-history-example,fig:superdtl:2-steps-ct-ex,fig:superdtl:crispr-scenario,fig:synesth:example,fig:synesth:crispr-sol,fig:adp-rec:rec-input,fig:adp-rec:syn-solution-best}.
To do so, the events of the reconciled tree must be laid out along the host tree while ensuring that sufficient space is reserved for each of them.
The program produces a {\LaTeX} document fragment that generates the diagram when integrated into a document.
Other programs exist for visualizing reconciliation scenarios (see \cite{Menet2022PhylogeneticReconciliation} for a recent review of the different options) but, as far as we know, none allows the visualization of segmental reconciliations.

When the scenario is time-consistent (see \cref{def:reconciliation:time-cons}), the program makes sure to position the events such that all horizontal transfers are represented by truly horizontal arcs.
Otherwise, the program attempts, through a greedy approach, to represent as many transfers as possible in this way and marks in red those for which it did not succeed.
This part of the program was written in collaboration with Arnaud~\textsc{Grandisson}.

The \textbf{\texttt{superrec2.compute}} module and the \texttt{\textbf{superrec2 reconcile}} command-line program make it possible to run the algorithm described in \cref{chap:adp-rec}.
This program makes it possible to apply different algebras for computing minimum-cost solutions or the set of Pareto-optimal vectors.
It can also enumerate the solutions or sample them randomly, either uniformly or according to a Boltzmann distribution whose temperature can be parameterized.
Statistics on the solution space can also be produced.

The software relies on the \emph{sowing} library for representing the manipulated trees and on the \emph{padapto} library for implementing the algorithms under the algebraic approach.
